\documentclass[11pt]{article}

\usepackage[margin=1in]{geometry}
\usepackage[T1]{fontenc}
\usepackage[utf8]{inputenc}
\usepackage{amsmath}
\usepackage{amssymb}
\usepackage{booktabs}
\usepackage{array}
\usepackage{enumitem}
\usepackage[colorlinks=true,linkcolor=blue,urlcolor=blue]{hyperref}

\setlength{\parindent}{0pt}
\setlength{\parskip}{0.55em}

\newcommand{\pts}[1]{\hspace{0.35em}\textnormal{\textbf{[#1 balai]}}}
\newenvironment{choices}{\begin{enumerate}[label=\Alph*.,leftmargin=2.2em,itemsep=0.12em,topsep=0.12em]}{\end{enumerate}}

\title{S181M116 Išvestinės finansinės priemonės ir alternatyvios investicijos\\
Baigiamasis egzaminas, B variantas: testas}
\author{Laikas: 1 valanda 30 minučių \qquad Iš viso: 80 balų}
\date{}

\begin{document}
\maketitle

\textbf{Instrukcijos.} Atsakykite į visus klausimus. Kiekviename punkte yra
vienas geriausias atsakymas. Jei nenurodyta kitaip, ignoruokite sandorių kaštus,
mokesčius, pirkimo ir pardavimo kainų skirtumus bei įsipareigojimų nevykdymo
riziką. Naudokite tolydinį palūkanų skaičiavimą. Piniginius atsakymus
apvalinkite iki dviejų skaitmenų po kablelio, tikimybes iki keturių skaitmenų po
kablelio, o apsidraudimo santykius iki keturių skaitmenų po kablelio.

\textbf{Vertinimas.} Egzamine yra 20 testo punktų. Kiekvienas punktas vertas 4
balų, iš viso 80 balų.

\textbf{Aprėptis.} Egzaminas remiasi klasėje aptarta medžiaga: išvestinių
priemonių pagrindais, rinkos struktūra, išankstiniais ir ateities sandoriais,
apsidraudimu ir bazės rizika, išankstinių ir ateities sandorių kainodara,
opcionais, pirkimo ir pardavimo opcionų paritetu, binominiais medžiais ir
opcionų nearbitražo apribojimais. 3 ir 5 paskaitų medžiaga neįtraukta.

\section*{1 klausimas: Sąvokos, išmokos ir rinkos struktūra \pts{16}}

Šis klausimas yra lengvesnis.

\begin{enumerate}
    \item Sutartis po šešių mėnesių išmoka $\max(90-S_T,0)$. Kuris teiginys
    teisingas? \pts{4}
    \begin{choices}
        \item Tai nėra išvestinė priemonė, nes nėra pradinės premijos.
        \item Tai išvestinė priemonė, kurios bazinis kintamasis yra turto kaina $S_T$.
        \item Tai obligacija, nes išmoka yra apribota iš viršaus.
        \item Tai ateities sandoris, nes išmoka yra tiesinė.
    \end{choices}

    \item Europos pirkimo opcionas ir Europos pardavimo opcionas turi tą pačią
    vykdymo kainą $K=75$. Pirkimo opciono premija yra 6, o pardavimo opciono
    premija yra 4. Jei $S_T=88$, kokios yra pirkėjo pirkimo opciono išmoka,
    pirkimo opciono pelnas, pardavimo opciono išmoka ir pardavimo opciono
    pelnas? \pts{4}
    \begin{choices}
        \item $13,\;7,\;0,\;-4$
        \item $13,\;13,\;0,\;0$
        \item $0,\;-6,\;13,\;9$
        \item $7,\;13,\;-4,\;0$
    \end{choices}

    \item Kuris teiginys geriausiai apibūdina OTC išvestinę priemonę, palyginti
    su biržoje prekiaujama išvestine priemone? \pts{4}
    \begin{choices}
        \item OTC sutartys yra labiau standartizuotos ir visada centralizuotai kliringuojamos.
        \item OTC sutartys yra pritaikomos, bet kelia daugiau sandorio šalies,
        vertinimo ir likvidumo rizikos.
        \item OTC sutartys neturi sandorio šalies rizikos, nes yra privačios.
        \item OTC sutartys visada pigesnės, nes joms nereikia garantinės įmokos.
    \end{choices}

    \item Kodėl garantinė įmoka gali sukurti likvidumo spaudimą? \pts{4}
    \begin{choices}
        \item Nes kasdienis atsiskaitymas gali priversti nuostolingas pozicijas
        įnešti pinigų dar negavus kompensuojančio bazinės pozicijos pinigų srauto.
        \item Nes garantinė įmoka yra ateities sandorio pirkimo kaina.
        \item Nes kintamoji garantinė įmoka pašalina visą kainos riziką.
        \item Nes garantinės įmokos reikalavimai atsiranda tik tada, kai
        apsidraudimas yra pelningas.
    \end{choices}
\end{enumerate}

\section*{2 klausimas: Ateities sandorių atsiskaitymas ir išankstinių sandorių kainodara \pts{16}}

Šis klausimas yra lengvesnis.

\begin{enumerate}
    \item Prekiautojas turi \textbf{trumpąją} poziciją aštuoniuose vario ateities
    sandoriuose. Kiekvienas sandoris yra 25 000 svarų. Vakar dienos atsiskaitymo
    kaina buvo 4,12 JAV dolerio už svarą, o šiandienos atsiskaitymo kaina yra
    4,05 JAV dolerio už svarą. Koks yra šiandienos pelnas arba nuostolis iš
    ateities sandorių? \pts{4}
    \begin{choices}
        \item 14 000 JAV dolerių pelnas
        \item 14 000 JAV dolerių nuostolis
        \item 1 400 JAV dolerių pelnas
        \item 1 400 JAV dolerių nuostolis
    \end{choices}

    \item Prieš atsiskaitymą prekiautojo garantinės sąskaitos likutis buvo
    38 000 JAV dolerių. Palaikomoji garantinė įmoka yra 30 000 JAV dolerių, o
    pradinė garantinė įmoka yra 38 000 JAV dolerių. Koks likutis po atsiskaitymo
    ir ar yra garantinės įmokos reikalavimas? \pts{4}
    \begin{choices}
        \item 24 000 JAV dolerių; taip, reikia įnešti 14 000 JAV dolerių.
        \item 30 000 JAV dolerių; garantinės įmokos reikalavimo nėra.
        \item 52 000 JAV dolerių; garantinės įmokos reikalavimo nėra.
        \item 52 000 JAV dolerių; taip, reikia įnešti 8 000 JAV dolerių.
    \end{choices}

    \item Akcijos dabartinė kaina yra $S_0=120$. Žinomų piniginių dividendų per
    šešių mėnesių išankstinio sandorio laikotarpį dabartinė vertė yra $I=2.40$.
    Tolydinė nerizikinga palūkanų norma yra 3,5\% per metus. Kokia yra
    nearbitražinė išankstinio sandorio kaina? \pts{4}
    \begin{choices}
        \item 117.60
        \item 119.68
        \item 122.10
        \item 124.25
    \end{choices}

    \item Prekiautojas tam pačiam sandoriui kotiruoja 121 išankstinę kainą.
    Kuris teiginys teisingas? \pts{4}
    \begin{choices}
        \item Kaina per maža; reikia parduoti turtą trumpai ir pirkti išankstinį sandorį.
        \item Kaina per didelė; reikia pirkti turtą, parduoti išankstinį sandorį
        ir užfiksuoti apie 1.32 pelno vienai akcijai.
        \item Kaina tiksliai atitinka nearbitražinę kainą.
        \item Kaina per didelė; reikia parduoti turtą trumpai, investuoti pajamas
        ir pirkti išankstinį sandorį.
    \end{choices}
\end{enumerate}

\section*{3 klausimas: Kryžminis apsidraudimas, bazės rizika ir efektyvios kainos \pts{16}}

Šis klausimas yra sunkesnis.

Gamintojas po trijų mėnesių pirks 1 250 000 svarų vario. Jis kryžmiškai
apsidraudžia naudodamas susijusio metalo ateities sandorį. Kiekvienas ateities
sandoris apima 25 000 svarų. Istoriniai duomenys yra:
\[
    \rho=0.72,\qquad \sigma_S=0.18,\qquad \sigma_F=0.16.
\]
Apsidraudimas pradedamas, kai $F_1=3.80$. Uždarant apsidraudimą, gamintojas
perka varį už $S_2=4.12$ ir uždaro ateities sandorius už $F_2=4.05$.

\begin{enumerate}
    \item Kokią apsidraudimo kryptį turėtų pasirinkti gamintojas? \pts{4}
    \begin{choices}
        \item Ilgąją ateities sandorių poziciją, nes jis vėliau pirks žaliavą ir
        nori apsisaugoti nuo kainos kilimo.
        \item Trumpąją ateities sandorių poziciją, nes jis vėliau pirks žaliavą.
        \item Ilgąją poziciją tik tada, jei laukiama bazė yra neigiama.
        \item Apsidraudimas neįmanomas, nes tai kryžminis apsidraudimas.
    \end{choices}

    \item Koks yra minimalios dispersijos apsidraudimo santykis? \pts{4}
    \begin{choices}
        \item 0.6400
        \item 0.7200
        \item 0.8100
        \item 1.1250
    \end{choices}

    \item Koks ateities sandorių skaičius gaunamas pagal apsidraudimo santykį
    prieš apvalinimą? \pts{4}
    \begin{choices}
        \item 36.0
        \item 40.5
        \item 50.0
        \item 56.25
    \end{choices}

    \item Kokia yra efektyvi pirkimo kaina vienam svarui ir kaip realizuota bazė
    palyginama su laukiama baze 0.03? \pts{4}
    \begin{choices}
        \item 3.87; realizuota bazė pirkėjui blogesnė.
        \item 3.87; realizuota bazė pirkėjui geresnė.
        \item 4.37; realizuota bazė pirkėjui blogesnė.
        \item 4.37; realizuota bazė pirkėjui geresnė.
    \end{choices}
\end{enumerate}

\section*{4 klausimas: Opcionų nearbitražas ir binominė kainodara \pts{16}}

Šis klausimas yra sunkesnis.

Nagrinėkite Europos opcionus akcijai, kuri nemoka dividendų. Dabartinė akcijos
kaina yra $S_0=48$, vykdymo kaina yra $K=50$, tolydinė nerizikinga palūkanų norma
yra $r=5\%$, o terminas yra $T=0.25$ metų. Pariteto klausimams Europos pirkimo
opcionas prekiaujamas už $c=2.60$. Atskirai, binominio modelio klausimams
naudokite $u=1.18$ ir $d=0.88$.

\begin{enumerate}
    \item Kokia nearbitražinė Europos pardavimo opciono kaina gaunama iš pirkimo
    ir pardavimo opcionų pariteto? \pts{4}
    \begin{choices}
        \item 1.22
        \item 2.60
        \item 3.20
        \item 3.98
    \end{choices}

    \item Tarkime, kad rinkos pardavimo opciono kaina yra $p=3.20$. Kuris
    arbitražo teiginys teisingas? \pts{4}
    \begin{choices}
        \item Apsauginis pardavimo opcionas brangus; parduoti jį ir pirkti
        fiduciarinį pirkimo opcioną.
        \item Fiduciarinis pirkimo opcionas brangus; parduoti jį ir pirkti
        apsauginį pardavimo opcioną.
        \item Abu paketai įkainoti teisingai.
        \item Arbitražo nėra, nes tiek pirkimo, tiek pardavimo opcionų kainos yra
        teigiamos.
    \end{choices}

    \item Vieno žingsnio medyje kokios yra akcijos kainos kilimo ir kritimo
    būsenose bei pirkimo opciono išmokos? \pts{4}
    \begin{choices}
        \item $S_u=56.64$, $S_d=42.24$; išmokos $6.64$ ir $0$.
        \item $S_u=56.64$, $S_d=42.24$; išmokos $0$ ir $7.76$.
        \item $S_u=59.00$, $S_d=44.00$; išmokos $9.00$ ir $0$.
        \item $S_u=48.18$, $S_d=47.88$; išmokos $0$ ir $0$.
    \end{choices}

    \item Kuri pora yra artimiausia rizikai neutraliai tikimybei ir binominei
    pirkimo opciono vertei? \pts{4}
    \begin{choices}
        \item $q=0.4419$, pirkimo opciono vertė $2.90$
        \item $q=0.4419$, pirkimo opciono vertė $6.64$
        \item $q=0.5581$, pirkimo opciono vertė $2.60$
        \item $q=0.5581$, pirkimo opciono vertė $3.67$
    \end{choices}
\end{enumerate}

\section*{5 klausimas: Opcionų ribos, ankstyvas vykdymas ir strategijų logika \pts{16}}

Šis klausimas yra sunkesnis.

Nagrinėkite opcionus akcijai, kuri nemoka dividendų. Dabartinė akcijos kaina yra
$S_0=70$, vykdymo kaina yra $K=68$, tolydinė nerizikinga palūkanų norma yra
$r=4.5\%$, o terminas yra $T=0.75$ metų.

\begin{enumerate}
    \item Kokia yra Europos pirkimo opciono apatinė riba ir ką ji reiškia, jei
    pirkimo opcionas kotiruojamas už $c=3.80$? \pts{4}
    \begin{choices}
        \item Apatinė riba apie 0; arbitražo nėra.
        \item Apatinė riba apie 1.74; arbitražo nėra.
        \item Apatinė riba apie 4.26; pirkimo opcionas pažeidžia apatinę ribą.
        \item Apatinė riba apie 6.00; pirkimo opcionas įkainotas teisingai.
    \end{choices}

    \item Kokia yra Europos pardavimo opciono apatinė riba ir ką ji reiškia, jei
    pardavimo opcionas kotiruojamas už $p=0.60$? \pts{4}
    \begin{choices}
        \item Apatinė riba 0; pagal šią ribą arbitražo nėra.
        \item Apatinė riba 1.74; pardavimo opcionas pažeidžia apatinę ribą.
        \item Apatinė riba 4.26; pardavimo opcionas pažeidžia apatinę ribą.
        \item Apatinė riba 68; pardavimo opcionas pažeidžia apatinę ribą.
    \end{choices}

    \item Jei Europos pirkimo opcionas su šia vykdymo kaina ir terminu
    prekiaujamas už $c=6.00$, kokia yra nearbitražinė Europos pardavimo opciono
    kaina pagal pirkimo ir pardavimo opcionų paritetą? \pts{4}
    \begin{choices}
        \item 0.00
        \item 1.74
        \item 3.80
        \item 6.00
    \end{choices}

    \item Kuris teiginys teisingas? \pts{4}
    \begin{choices}
        \item Amerikos pardavimo opciono kaina $P=1.50$ suderinama, jei Europos
        pardavimo opciono vertė pagal paritetą yra apie 1.74.
        \item Dengtas pirkimo opcionas yra nerizikingas, nes opciono išrašytojas turi akciją.
        \item Amerikos pardavimo opcionas turi būti vertas bent tiek, kiek
        atitinkamas Europos pardavimo opcionas, o dengtas pirkimo opcionas vis
        dar turi akcijos kainos kritimo riziką.
        \item Dengtas pirkimo opcionas turi neribotą augimo potencialą, nes
        opciono išrašytojas turi akciją.
    \end{choices}
\end{enumerate}

\end{document}
